% SPDX-License-Identifier: GPL-3.0-or-later
% Copyright (C) 2026 Alessandro Greco
\chapter{Gestire un relay o un bridge}\label{cap:22-relay}
\aggiornatoal{23 settembre 2026}{tor 0.4.9.13 (manuale); rete di laboratorio}

\begin{obiettivi}
Al termine di questo capitolo saprai:
\begin{itemize}[nosep]
  \item scegliere tra relay intermedio, exit e bridge in base alle tue risorse e ai rischi che puoi sostenere;
  \item scrivere la configurazione essenziale di ciascuno, con limiti di banda e chiavi custodite offline;
  \item leggere e scrivere una politica di uscita, e verificarne l'effetto nella rete di laboratorio (Laboratorio~14);
  \item dichiarare correttamente una famiglia di relay.
\end{itemize}
\end{obiettivi}

La rete Tor esiste perché migliaia di persone e organizzazioni mettono a disposizione banda e computer. Nei capitoli precedenti abbiamo guardato i relay dal punto di vista del client; qui passiamo dall'altra parte. Non è un manuale di amministrazione, che il Tor Project mantiene aggiornato nella propria guida per gli operatori, ma una spiegazione di che cosa significano le scelte principali alla luce del protocollo.

% ---------------------------------------------------------------------------
\section{Che tipo di nodo}\label{sec:relay-tipi}

\begin{description}
  \item[Relay intermedio] (\emph{middle}, e col tempo eventualmente guard): inoltra celle tra altri relay tramite connessioni Internet, ma non apre connessioni verso i siti di destinazione dei client. È il ruolo meno impegnativo: il suo indirizzo compare nel consenso, ma nessun sito vede traffico provenire da esso.
  \item[Exit] apre, per conto dei client, le connessioni verso le destinazioni. È il ruolo più prezioso per la rete, perché gli exit sono pochi rispetto alla domanda (\cref{sec:percorso-pesi}), ed è quello più esposto: per chi riceve un abuso, l'autore apparente è l'indirizzo dell'exit.
  \item[Bridge] è un relay non elencato nel consenso, usato come primo salto da chi non può raggiungere la rete pubblica (\cref{sec:cens-bridge}). Con un pluggable transport come obfs4 o WebTunnel, anche un collegamento domestico modesto è utile.
\end{description}

\begin{attenzione}[Prima di cominciare]
Gestire un exit ha conseguenze legali e pratiche (segnalazioni di abuso, richieste delle autorità) che discutiamo nel \cref{cap:25-legale}. Conviene iniziare da un relay intermedio o da un bridge, informare il fornitore di connettività, e non usare mai per un relay la stessa connessione da cui si usa Tor come client per attività sensibili.
\end{attenzione}

% ---------------------------------------------------------------------------
\section{La configurazione essenziale}\label{sec:relay-config}

Un relay intermedio richiede poche righe di \file{torrc}:
\begin{lstlisting}[numbers=none]
Nickname        MioRelay
ContactInfo     operatore@example.org
ORPort          443
ExitRelay       0
SocksPort       0
RelayBandwidthRate  4 MBytes
RelayBandwidthBurst 8 MBytes
\end{lstlisting}
\opt{ContactInfo} serve a essere contattati se il relay ha problemi; il manuale ricorda che i descrittori sono archiviati e pubblicati, quindi l'indirizzo sarà pubblico per sempre. \opt{ExitRelay 0} rende esplicito ciò che nelle versioni recenti è già il comportamento predefinito: con il valore \texttt{auto}, tor fa da exit solo se l'operatore ha scritto una politica di uscita. \opt{RelayBandwidthRate} e \opt{RelayBandwidthBurst} limitano la banda media e di picco; \opt{AccountingMax} fissa invece un tetto al traffico di un periodo (per esempio un mese), raggiunto il quale il relay smette di accettare nuovo traffico fino al periodo successivo.

\textbf{Nota operativa.} Con il valore automatico di \texttt{ExitRelay}, anche \texttt{ReducedExitPolicy} o \texttt{IPv6Exit} possono abilitare il ruolo di uscita: non conta soltanto una riga \texttt{ExitPolicy}. La direttiva esplicita \texttt{ExitRelay 0} dell'esempio impedisce invece qualunque uscita. Verifica sempre il risultato prima di pubblicare un relay reale.

\subsection{La vita di un relay nuovo}

Un relay appena avviato pubblica il proprio descrittore alle autorità, che lo inseriscono nel consenso e iniziano a misurarlo (\cref{sec:dir-voto}). Nei primi giorni riceve pochissimo traffico: i client scelgono i relay in proporzione ai pesi di banda, e finché le misure non sono stabili il peso resta basso. I flag arrivano col tempo (\cref{sec:dir-flag}): \texttt{Stable} richiede che il relay resti acceso a lungo senza interruzioni, e \texttt{Guard} richiede anche un'anzianità minima e un'alta disponibilità. Diventato guard, il relay vedrà prima calare il traffico, perché i client che già hanno i propri guard non lo scelgono come intermedio con la stessa frequenza, e poi crescere lentamente, man mano che nuovi client lo adottano (\cref{sec:guard}). I due casi che fanno perdere terreno non sono equivalenti: chi riavvia spesso il relay, senza cambiarne l'identità, conserva le misure di banda ma perde la disponibilità continua richiesta per i flag \texttt{Stable} e \texttt{Guard}, e deve riconquistarli; chi ne cambia le \emph{chiavi d'identità} crea invece un relay del tutto nuovo per le autorità, che riparte da zero anche nelle misure di banda.

\subsection{Le chiavi}

L'identità di un relay è la sua chiave ed25519 principale (\cref{sec:canale-chiavi}). Per default tor la tiene sul disco del server; con \opt{OfflineMasterKey 1} la chiave principale resta su un altro computer, e sul server vanno solo chiavi di firma temporanee, certificate con \cmd{tor --keygen}. Se il server viene compromesso, l'attaccante ottiene chiavi che scadono, non l'identità del relay. È un accorgimento che richiede disciplina: bisogna rinnovare le chiavi temporanee prima che scadano, altrimenti il relay smette di funzionare.

% ---------------------------------------------------------------------------
\section{La politica di uscita}\label{sec:relay-exitpolicy}

La politica di uscita è una lista di regole \texttt{accept} e \texttt{reject}, ciascuna su un indirizzo (o una rete) e una porta (o un intervallo). Tor le esamina \textbf{nell'ordine} e applica la prima che corrisponde alla destinazione. Le regole scritte dall'operatore precedono quelle predefinite, che il manuale elenca: vengono rifiutate le porte della posta (25), di NNTP (119), di NetBIOS e SMB (135--139, 445), di alcuni programmi di condivisione file, e poi tutto il resto è accettato. Indipendentemente da questo, tor rifiuta per default le destinazioni private (\texttt{127.0.0.0/8}, \texttt{10.0.0.0/8}, \texttt{192.168.0.0/16} e simili) e il proprio indirizzo pubblico, per non diventare un accesso alla rete interna dell'operatore (opzione \opt{ExitPolicyRejectPrivate}, che la rete di laboratorio disattiva per poter usare \texttt{127.0.0.1}).

\textbf{Ordine delle protezioni.} Quando \texttt{ExitPolicyRejectPrivate} è attivo (valore predefinito), i rifiuti verso reti private e indirizzi locali vengono anteposti alle regole definite dall'operatore, non semplicemente aggiunti in fondo. La rete Chutney dell'esperimento disabilita esplicitamente questo controllo: non replicare questa impostazione su un exit pubblico.

Un exit che accetta solo il web si scrive così:
\begin{lstlisting}[numbers=none]
ExitRelay  1
ExitPolicy accept *:80
ExitPolicy accept *:443
ExitPolicy reject *:*
\end{lstlisting}
Il Tor Project suggerisce, come compromesso diffuso, la \opt{ReducedExitPolicy}, che accetta un elenco di porte di servizi comuni ed esclude quelle che generano più segnalazioni.

\begin{figure}[htbp]
\centering
% SPDX-License-Identifier: GPL-3.0-or-later
% Copyright (C) 2026 Alessandro Greco
%% Testo alternativo per l'accessibilita' (usato solo nell'EPUB).
\altfig{Politica di uscita dell'exit del laboratorio Chutney: ExitPolicyRejectPrivate 0 disabilita soltanto nella rete di test il rifiuto preventivo delle destinazioni private. Due richieste BEGIN verso 127.0.0.1:8080 e 127.0.0.1:8081 attraversano le regole in ordine. La porta 8080 corrisponde prima a reject 127.0.0.1:8080 e riceve END EXITPOLICY. La 8081 non corrisponde al reject e passa grazie a accept 127.0.0.0/8:*, che consente di tentare una connessione TCP; CONNECTED richiede un server in ascolto. Se accept precedesse reject, il rifiuto sarebbe inefficace. Non disabilitare ExitPolicyRejectPrivate su exit pubblici.}
% Politica di uscita nel solo laboratorio
\begin{tikzpicture}[font=\sffamily\scriptsize,node distance=0.45cm,
  reg/.style={draw,rounded corners=1pt,minimum height=0.6cm,text width=4.6cm,font=\ttfamily\scriptsize},
  esito/.style={draw,rounded corners=2pt,minimum height=0.6cm,text width=2.5cm,align=center}]
  \node[font=\sffamily\scriptsize] (top) at (0,1.25) {due BEGIN: 127.0.0.1:8080 / :8081};
  \node[draw,rounded corners=2pt,fill=torviola!10,font=\ttfamily\tiny] (r0) at (0,0.15) {ExitPolicyRejectPrivate 0};
  \node[reg,fill=avviso!12,below=of r0] (r1) {reject 127.0.0.1:8080};
  \node[reg,fill=torverde!12,below=of r1] (r2) {accept 127.0.0.0/8:*};
  \node[reg,fill=torgrigio!10,below=of r2] (r3) {(altre regole predefinite)};
  \draw[->] (top) -- (r0); \draw[->] (r0) -- (r1); \draw[->] (r1) -- node[right,font=\sffamily\tiny] {8081: non corrisponde} (r2); \draw[->] (r2) -- node[right,font=\sffamily\tiny] {altrimenti} (r3);
  \node[esito,fill=avviso!15,draw=avviso,right=1.2cm of r1] (e1) {\texttt{END} motivo\\\texttt{EXITPOLICY}};
  \draw[->,avviso,thick] (r1) -- node[above,font=\sffamily\tiny] {8080 corrisponde} (e1);
  \node[esito,fill=torverde!15,draw=torverde!60!black,right=1.2cm of r2] (e2) {tentativo TCP\\CONNECTED solo\\se server attivo};
  \draw[->,torverde!60!black,dashed] (r2) -- node[above,font=\sffamily\tiny] {8081 corrisponde} (e2);
  \node[anchor=west,font=\sffamily\tiny,text=torgrigio,align=left] at (-5.9,-0.55) {vince la prima regola\\che corrisponde};
  \node[anchor=west,font=\sffamily\tiny,text=avviso] at (-5.9,-1.95) {non su exit pubblico};
  \node[font=\sffamily\tiny,text=torgrigio,anchor=south] at (r0.north) {solo laboratorio: ExitPolicyRejectPrivate 0};
\end{tikzpicture}

\caption{Esempio del laboratorio Chutney: \opt{ExitPolicyRejectPrivate 0} è impostata soltanto per consentire indirizzi privati nella rete di prova; normalmente Tor antepone il loro rifiuto alle regole dell'operatore. La richiesta \texttt{BEGIN} verso \texttt{127.0.0.1:8080} corrisponde a \texttt{reject} e produce \texttt{END EXITPOLICY}; la porta 8081 supera tale regola e viene accettata. \texttt{accept} consente il tentativo TCP, non garantisce \texttt{CONNECTED} senza un server in ascolto. Invertire le regole renderebbe inefficace il rifiuto.}
\label{fig:policy}
\end{figure}

\paragraph{Che cosa arriva ai client.} Il descrittore del relay contiene la politica completa; il consenso e i microdescrittori ne portano un riassunto, solo per porte (\texttt{p accept 80,443}), da cui il client decide se un exit può servire una richiesta (\cref{sec:percorso-vincoli}). Il riassunto non contiene indirizzi: se l'exit rifiuta una destinazione per indirizzo, il client lo scopre solo provando, come mostra il laboratorio.

% ---------------------------------------------------------------------------
\section{Laboratorio: una politica di uscita alla prova}\label{sec:lab-relay}

\begin{laboratorio}[14]{Una regola di uscita e le sue conseguenze}
\textbf{Obiettivo.} Aggiungere una regola a un exit della rete di laboratorio e osservare, con \file{labs/it/codice/torlab}, la risposta a destinazioni accettate e rifiutate.\\
\textbf{Prerequisiti.} Laboratori~3 e~4; la rete di laboratorio avviata.\\
\textbf{Tempo stimato.} 1 ora.
\end{laboratorio}

\paragraph{Passo 1: la regola.} Nel \file{torrc} di \texttt{test006e} aggiungiamo, prima delle regole esistenti che accettano la rete \texttt{127.0.0.0/8}, una riga che rifiuta una sola porta:
\begin{lstlisting}[numbers=none]
ExitPolicy reject 127.0.0.1:8080
ExitPolicy accept 127.0.0.0/8:*
\end{lstlisting}
Il listato mostra l'ordine delle due regole interessate: quando modifichi il \file{torrc} già esistente, aggiungi soltanto la nuova riga \texttt{reject} prima della prima \texttt{ExitPolicy}, senza duplicare la regola \texttt{accept}. L'ordine conta: messa dopo la regola \texttt{accept}, la nostra riga non avrebbe alcun effetto, perché la prima regola che corrisponde vince. È esattamente l'errore registrato nel precedente esperimento di riferimento. Dopo aver riavviato il relay, avviamo due piccoli server HTTP sulle porte 8080 e 8081 del container.

\paragraph{Passo 2: tre richieste.} Con \texttt{costruisci} e il modulo \file{stream.py} apriamo un circuito che termina su \texttt{006e} e inviamo tre \texttt{BEGIN}, verso le porte 8080, 8081 e 9999 (su quest'ultima non ascolta nessuno). Le risposte, in una nostra sessione (codici, indirizzo e TTL possono variare secondo la versione e l'ambiente):
\begin{lstlisting}[numbers=none]
127.0.0.1:8080 -> END  motivo=EXITPOLICY indirizzo=127.0.0.1 TTL=1800
127.0.0.1:8081 -> CONNECTED
127.0.0.1:9999 -> END  motivo=MISC
\end{lstlisting}

\paragraph{Passo 3: leggere le risposte.} Il rifiuto per politica è comunicato con il motivo \texttt{EXITPOLICY}, e la cella \texttt{END} porta con sé l'indirizzo risolto e un TTL: così il client può ricordare che quell'exit non serve quella destinazione e sceglierne un altro, senza aspettare un nuovo descrittore. La porta senza server, invece, produce un generico \texttt{MISC}, e non il \texttt{CONNECTREFUSED} che la specifica elenca per le connessioni rifiutate \spec{tor-spec}{Closing streams}.

\begin{approfondimento}[Perché \texttt{MISC} e non \texttt{CONNECTREFUSED}]
Il motivo si trova nei log dell'exit, al livello \texttt{info}, e vale la pena seguirlo perché mostra come si indaga una differenza tra specifica e implementazione:
\begin{lstlisting}[numbers=none]
[info] connection_handle_write_impl(): in-progress connect failed.
       Removing. (Connection refused)
[info] errno_to_stream_end_reason(): Didn't recognize errno 115 (Operation
       now in progress); telling the client that we are ending a stream
       for 'misc' reason.
\end{lstlisting}
L'exit \emph{sa} che la connessione è stata rifiutata, ma sceglie il motivo a partire da un altro valore. La connessione verso la destinazione non è bloccante: \texttt{connect()} ritorna subito con \texttt{EINPROGRESS} (errno 115) e l'esito arriva più tardi, come evento di scrittura. A quel punto \file{connection.c} legge l'errore vero con \texttt{getsockopt(SO\_ERROR)}, lo registra nel log (\enquote{Connection refused}) e chiama \file{connection\_edge\_end\_errno()}. Questa funzione, però, ricalcola il motivo con \texttt{errno\_to\_stream\_end\_reason(tor\_socket\_errno(\ldots))}: su Unix \texttt{tor\_socket\_errno} è definito come la macro \texttt{(errno)} (\file{src/lib/net/socket.h}), cioè la variabile globale, che conserva ancora \texttt{EINPROGRESS}. Il valore 115 non compare nell'elenco di \file{src/core/or/reasons.c}, e il ramo predefinito restituisce \texttt{MISC}. L'informazione giusta esiste, ma non è quella che viene usata. Nello stesso blocco di codice le due diramazioni non concordano: quella per le connessioni tra relay passa il valore letto, quella per gli stream lo ignora. Secondo il resoconto e la patch già presenti in \file{docs/it/segnalazioni/}, passando il valore corretto anche lì, nella precedente prova privata su tor 0.4.9.13 la terza richiesta ha prodotto \texttt{CONNECTREFUSED}. È un difetto minore, senza conseguenze per l'anonimato: il client sa comunque che lo stream è fallito, ma perde un'informazione su cui agirebbe, perché \texttt{CONNECTREFUSED} gli direbbe che l'exit ha raggiunto la destinazione e che riprovare altrove è inutile. Analisi e patch sono in \file{docs/it/segnalazioni/}, pronte per il Tor Project. È anche un buon esempio di perché conviene verificare sul vivo ciò che si legge nelle specifiche.
\end{approfondimento}

Lo script \file{labs/it/laboratorio/politica\_uscita.py}, eseguibile dall'host Linux dalla radice della repository con \file{PYTHONPATH=labs/it/codice python3 labs/it/laboratorio/politica\_uscita.py}, esegue i tre tentativi e tenta di ripristinare il \file{torrc} originale di \texttt{006e} tramite un blocco \texttt{finally}. Usalo soltanto nella rete Chutney dedicata: modifica temporaneamente la configurazione dell'exit e apre server HTTP locali. Fa parte di \texttt{make verifica-lab}; il target \texttt{make verifica} della cartella del laboratorio, invece, controlla soltanto la rete Chutney. Un dettaglio necessario: il circuito ha tre salti, perché un exit rifiuta di far uscire il traffico da un circuito di un solo salto e risponderebbe \texttt{END} con motivo \texttt{TORPROTOCOL} (\cref{sec:dos-limiti}).

% ---------------------------------------------------------------------------
\section{Le famiglie}\label{sec:relay-famiglie}

Chi gestisce più relay deve dichiararli come \textbf{famiglia}: il client non usa mai due relay della stessa famiglia nello stesso circuito (\cref{sec:percorso-vincoli}), perché chi li controlla entrambi potrebbe correlare il traffico. Il meccanismo storico è \opt{MyFamily}: ogni relay elenca le impronte di tutti gli altri, e la famiglia vale solo se le dichiarazioni sono reciproche. Con dieci relay sono dieci elenchi da tenere allineati, e ogni aggiunta richiede di modificarli tutti.

Il meccanismo introdotto con tor 0.4.9 usa invece una \textbf{chiave di famiglia}: l'operatore la genera una volta con \cmd{tor --keygen-family MyFamilyKey}, copia il file risultante \texttt{MyFamilyKey.secret\_family\_key} nella directory delle chiavi (\texttt{KeyDirectory}) di ogni relay, mantenendo il suffisso \texttt{.secret\_family\_key}, e inserisce in ciascun \file{torrc} la riga \texttt{FamilyId} generata dallo strumento. Proteggi e conserva questa chiave segreta: se viene compromessa, va ruotata sull'intera famiglia. Ogni relay firma con quella chiave un certificato che lo lega alla famiglia, e chi verifica non ha più bisogno di confrontare elenchi. Nella fase di transizione alcune versioni hanno supportato entrambi i meccanismi. Per nuove installazioni occorre seguire la guida aggiornata del Tor Project: FamilyId e la chiave condivisa sono il metodo corrente, mentre MyFamily è stato rimosso nelle versioni più recenti. Il meccanismo precedente resta descritto qui per interpretare le configurazioni storiche.

\begin{approfondimento}[Perché dichiararsi conviene]
Nessuno obbliga un operatore a dichiarare la famiglia, e un avversario non lo farebbe. Ma le autorità e i volontari che sorvegliano la rete cercano relay con caratteristiche comuni (stessa rete, stessa versione, avvio simultaneo, stesso contatto): una famiglia non dichiarata viene notata e può portare all'esclusione dei relay dal consenso. Dichiararla è parte del patto di fiducia su cui si regge la rete.
\end{approfondimento}

% ---------------------------------------------------------------------------
\textbf{I bridge non appartengono agli elenchi pubblici \texttt{MyFamily}:} associarli ai relay pubblici potrebbe comprometterne la riservatezza. La guida di Tor raccomanda di non usare questa opzione sui bridge. Inoltre \texttt{tor --keygen-family MyFamilyKey} può sovrascrivere file esistenti con quel nome: usa un nome nuovo e custodisci il file segreto con permessi stretti.

\section{Un bridge}\label{sec:relay-bridge}

Un bridge obfs4 si configura come nel Laboratorio~10 (\cref{sec:lab-censura}):
\begin{lstlisting}[numbers=none]
BridgeRelay 1
ORPort 9001
ServerTransportPlugin obfs4 exec /usr/bin/obfs4proxy
ServerTransportListenAddr obfs4 0.0.0.0:443
ExtORPort auto
BridgeDistribution any
\end{lstlisting}
\opt{BridgeDistribution} indica al servizio di distribuzione (oggi rdsys) come vorremmo che il bridge fosse dato ai client: \texttt{any} lascia decidere al servizio, \texttt{none} lo esclude del tutto, utile per un bridge privato da dare a mano a persone di fiducia; i canali disponibili (sito web, posta, Telegram e altri) sono elencati nella pagina informativa dei bridge. Un bridge non compare nel consenso, ma il suo descrittore va comunque all'autorità dei bridge, che ne verifica la raggiungibilità.

Conviene che la porta del transport sia diversa dall'ORPort e, se possibile, una porta comune come 443: è quella che i client usano e che un censore vede (\cref{sec:cens-strategie}).

% ---------------------------------------------------------------------------
\section*{Riepilogo}
\begin{itemize}
  \item Un relay intermedio è il modo meno impegnativo per contribuire; un exit è il più utile e il più esposto; un bridge aiuta chi è censurato.
  \item Poche opzioni bastano: identità, contatto, porte, limiti di banda; la chiave principale può restare offline.
  \item La politica di uscita è una lista ordinata in cui vince la prima regola che corrisponde; le destinazioni private sono rifiutate per default.
  \item Un rifiuto per politica arriva al client come \texttt{END} con motivo \texttt{EXITPOLICY}, indirizzo e TTL.
  \item Più relay dello stesso operatore vanno dichiarati come famiglia, con \opt{MyFamily} o con la nuova chiave di famiglia.
\end{itemize}

% ---------------------------------------------------------------------------
\section*{Esercizi}

\begin{esercizio}[1]\label{es:22-ordine}
Data la politica \texttt{accept *:80}, \texttt{reject 192.0.2.0/24:*}, \texttt{reject *:*}, un client può uscire verso \texttt{192.0.2.7:80}? Riscrivi la politica perché non possa.
\end{esercizio}

\begin{esercizio}[1]\label{es:22-riassunto}
Perché il riassunto per porte nel consenso non basta a sapere se un exit accetterà una connessione? Che cosa fa il client quando riceve \texttt{END} con motivo \texttt{EXITPOLICY}?
\end{esercizio}

\begin{esercizio}[2]\label{es:22-famiglia}
Un operatore gestisce cinque relay con \opt{MyFamily} e ne aggiunge un sesto, ma aggiorna la configurazione solo sul nuovo. Che cosa vedono i client? E con la chiave di famiglia?
\end{esercizio}

\begin{esercizio}[2]\label{es:22-refused}
Verifica la spiegazione dell'approfondimento: attiva sull'exit i log di livello \texttt{info} e ripeti le tre richieste. Che cosa cambierebbe se \file{connection\_edge\_end\_errno()} ricevesse l'errore letto con \texttt{getsockopt}? Quali altri motivi di chiusura sarebbero riportati correttamente?
\end{esercizio}

\begin{esercizio}[3]\label{es:22-accounting}
Un relay ha \texttt{AccountingMax 1 TBytes} su base mensile e una banda di 10\,MB/s. Dopo quanti giorni, a pieno regime, esaurisce il tetto? Quali effetti ha sulla rete un relay che resta spento per metà del mese, rispetto a uno che funziona tutto il mese con banda dimezzata?
\end{esercizio}

% ---------------------------------------------------------------------------
\section*{Per approfondire}
\begin{itemize}
  \item La guida per gli operatori di relay del Tor Project, \url{https://community.torproject.org/relay/}.
  \item \begin{sloppypar}Il manuale di tor, alle voci \opt{ExitPolicy}, \opt{ReducedExitPolicy},
        \opt{FamilyId}, \opt{OfflineMasterKey} e \opt{BridgeDistribution}.\end{sloppypar}
  \item La specifica dei motivi di chiusura degli stream, \spec{tor-spec}{Closing streams}.
\end{itemize}
